% Compile with: xelatex report.tex (run twice for the table of contents)
\documentclass[11pt,a4paper]{article}

\usepackage[a4paper,top=2.4cm,bottom=2.4cm,left=2.3cm,right=2.3cm]{geometry}
\usepackage{fontspec}
\usepackage{amsmath}
\usepackage{graphicx}
\usepackage{array}
\usepackage{booktabs}
\usepackage{tabularx}
\usepackage{float}
\usepackage{xcolor}
\usepackage{tikz}
\usetikzlibrary{arrows.meta}
\usepackage{fancyvrb}
\usepackage{needspace}
\usepackage{titlesec}
\usepackage{fancyhdr}
\usepackage[hidelinks]{hyperref}
\usepackage{xepersian}

\settextfont[Scale=1.1]{Times New Roman}
\setlatintextfont{Times New Roman}
\setlatinmonofont[Scale=0.88]{Consolas}
\defpersianfont\headfont[Scale=0.98]{Tahoma}

\definecolor{accent}{HTML}{1F4E79}
\definecolor{softbg}{HTML}{EEF3F8}

\titleformat*{\section}{\large\bfseries\headfont\color{accent}}
\titleformat*{\subsection}{\normalsize\bfseries\headfont\color{accent}}
\titlespacing*{\section}{0pt}{16pt}{7pt}
\titlespacing*{\subsection}{0pt}{11pt}{5pt}

\setlength{\emergencystretch}{3em}
\setlength{\parindent}{0pt}
\setlength{\parskip}{6pt}
\linespread{1.32}
\renewcommand{\arraystretch}{1.28}

\pagestyle{fancy}
\fancyhf{}
\fancyhead[C]{\small\color{gray} گزارش نهایی --- سامانه پردازش گفتار فارسی و شبیه‌سازی صدا}
\fancyfoot[C]{\small\thepage}
\renewcommand{\headrulewidth}{0.3pt}

\newcommand{\en}[1]{\lr{#1}}
\newcommand{\code}[1]{\lr{\texttt{#1}}}
\newcolumntype{R}[1]{>{\raggedleft\arraybackslash}p{#1}}
\newcolumntype{L}[1]{>{\raggedright\arraybackslash}p{#1}}
\DefineVerbatimEnvironment{codeblock}{Verbatim}{fontsize=\footnotesize,frame=single,rulecolor=\color{gray!50},framesep=6pt,baselinestretch=1.05}

\begin{document}

\begin{center}
  {\headfont\bfseries\LARGE\color{accent} سامانه پردازش گفتار فارسی و شبیه‌سازی صدا\par}
  \vspace{4pt}
  {\large گزارش نهایی سؤال دوم، ارزیابی عملی هوش مصنوعی\par}
  \vspace{4pt}
  {\large \en{Mohammad Mahdi Abedi}\par}
\end{center}
\vspace{6pt}

\noindent\colorbox{softbg}{\parbox{\dimexpr\linewidth-2\fboxsep}{%
\textbf{خلاصه.} برای سؤال دوم سامانه‌ای ساختم که فایل صوتی فارسی را به متن تبدیل می‌کند و با چند ثانیه از صدای یک گوینده، متن جدیدی را با صدای مشابه او می‌خواند. هر دو بخش از یک رابط وب در دسترس‌اند و مدل‌ها روی سرور خودمان اجرا می‌شوند.

مدل گفتار به متن روی ۳۰۰ جمله آزمون \en{WER} برابر ۸٫۸۸٪ (بازه اطمینان ۹۵٪: ۷٫۸۷٪ تا ۹٫۹۵٪) و \en{CER} برابر ۲٫۸٪ گرفت؛ \en{WER} مدل پایه روی همین داده ۱۷٫۹۶٪ بود. در شبیه‌سازی صدا، میانگین شباهت گوینده روی ۱۲ نمونه ۰٫۶۹ و \en{CER} رفت‌وبرگشتی ۳٫۴٪ شد. طبیعی‌بودن صدا را با عدد نسنجیده‌ام و گفتار محاوره‌ای را هم آزمون نکرده‌ام؛ جزئیات این دو در بخش~\ref{sec:limits} آمده است.
}}

\vspace{4pt}
\setcounter{tocdepth}{1}
\tableofcontents
\clearpage

\section{مقدمه و دامنه کار}

سؤال دوم دو کار می‌خواهد: تبدیل یک فایل صوتی فارسی ۳۰ تا ۶۰ ثانیه‌ای به متن، و خواندن یک پاراگراف جدید با صدایی شبیه نمونه ورودی، همه در یک رابط ساده.

چون صدای افراد داده حساسی است، تصمیم گرفتم همه چیز روی سرور خودمان اجرا شود و صوتی به سرویس بیرونی نرود. در ده روز فرصت، آموزش مدل از صفر منطقی نبود؛ پس از مدل‌های آماده‌ای شروع کردم که برای فارسی تنظیم شده‌اند و بیشتر وقتم را روی بررسی خروجی‌ها و رفع خطاهایی گذاشتم که واقعاً می‌دیدم. نتایج اصلی را می‌شود با اسکریپت‌های مخزن دوباره تولید کرد (بخش~\ref{sec:guide}).

\section{الزامات و وضعیت پاسخ}

\begin{table}[H]
\centering\small
\begin{tabularx}{\linewidth}{R{4.6cm} X R{3.1cm}}
\toprule
\textbf{الزام} & \textbf{پیاده‌سازی} & \textbf{وضعیت} \\
\midrule
\multicolumn{3}{c}{\textbf{بخش اول: گفتار به متن}} \\
\midrule
پشتیبانی از فایل صوتی فارسی & خواندن قالب‌های رایج با \en{ffmpeg} و تبدیل به ۱۶ کیلوهرتز تک‌کانال & انجام شده \\
کمترین خطای ممکن & \en{Whisper large-v3-turbo} تنظیم‌شده برای فارسی و برش صوت در مکث & \en{WER} برابر ۸٫۹٪ \\
تلفظ‌ها و گفتار محاوره‌ای & فقط از راه انتخاب مدل؛ سازوکار اختصاصی ساخته نشده است & آزمون نشده \\
سنجش با متن مرجع & محاسبه \en{WER} و \en{CER} در رابط و در اسکریپت ارزیابی & انجام شده \\
\midrule
\multicolumn{3}{c}{\textbf{بخش دوم: متن به گفتار با صدای گوینده}} \\
\midrule
دریافت نمونه صدا و متن فارسی & زبانه دوم رابط & انجام شده \\
تولید صدای مشابه نمونه & \en{XTTS-v2} تنظیم‌شده روی پیکره \en{ParsVoice} & شباهت ۰٫۶۹ \\
تلفظ، روانی و طبیعی‌بودن & آماده‌سازی متن و تنظیم پارامترها؛ وضوح با \en{CER} رفت‌وبرگشتی سنجیده شده است & وضوح سنجیده شده؛ طبیعی‌بودن نیازمند شنیدن \\
متن‌های خارج از نمونه صدا & سه پاراگراف آزمون که در نمونه‌های صدا وجود ندارند & انجام شده \\
رضایت صاحب صدا & چک‌باکس اجباری در رابط & انجام شده \\
\bottomrule
\end{tabularx}
\caption{الزامات سؤال و وضعیت هر کدام.}
\end{table}

\section{معماری سامانه}\label{sec:arch}

\begin{figure}[H]
\centering
\begin{latin}
\begin{tikzpicture}[
  box/.style={draw=accent,fill=softbg,rounded corners=2pt,minimum height=15mm,text width=24mm,align=center,font=\footnotesize},
  arr/.style={-{Stealth[length=2.2mm]},thick,accent}]
  \node[box] (a1) at (13.6,0) {\rl{فایل صوتی فارسی}};
  \node[box] (a2) at (10.2,0) {\lr{ffmpeg}\\\rl{۱۶ کیلوهرتز}\\\rl{تک‌کانال}};
  \node[box] (a3) at (6.8,0) {\rl{برش در مکث}\\\rl{حداکثر ۲۸ ثانیه}};
  \node[box] (a4) at (3.4,0) {\lr{Whisper turbo}\\\rl{فارسی}};
  \node[box] (a5) at (0,0) {\rl{متن}\\\lr{WER / CER}};
  \draw[arr] (a1) -- (a2); \draw[arr] (a2) -- (a3); \draw[arr] (a3) -- (a4); \draw[arr] (a4) -- (a5);

  \node[box] (b1) at (13.6,-2.5) {\rl{نمونه صدا}\\\rl{و متن جدید}};
  \node[box] (b2) at (10.2,-2.5) {\rl{یکسان‌سازی بلندی}\\\rl{و بردار گوینده}};
  \node[box] (b3) at (6.8,-2.5) {\rl{عدد به حروف،}\\\rl{جمله‌بندی}};
  \node[box] (b4) at (3.4,-2.5) {\lr{XTTS-v2}\\\lr{ParsVoice}\\\rl{جمله‌به‌جمله}};
  \node[box] (b5) at (0,-2.5) {\rl{صدای ۲۴ کیلوهرتز}\\\rl{شباهت \lr{ECAPA}}};
  \draw[arr] (b1) -- (b2); \draw[arr] (b2) -- (b3); \draw[arr] (b3) -- (b4); \draw[arr] (b4) -- (b5);
\end{tikzpicture}
\end{latin}
\caption{مسیر پردازش. ردیف بالا گفتار به متن و ردیف پایین شبیه‌سازی صدا است.}
\end{figure}

\subsection{اجزای نرم‌افزار}

\begin{table}[H]
\centering\small
\begin{tabularx}{\linewidth}{L{5.6cm} X}
\toprule
\textbf{فایل} & \textbf{کار} \\
\midrule
\code{app.py} & رابط کاربری \en{Gradio} با دو زبانه \\
\code{pv/audio.py} & خواندن هر قالب صوتی با \en{ffmpeg}، ذخیره \en{wav} و یکسان‌سازی بلندی \\
\code{pv/text.py} & یکسان‌سازی متن برای امتیازدهی و برای تولید گفتار، عدد به حروف، جمله‌بندی \\
\code{pv/stt.py} & برش در مکث و رونویسی با \en{Whisper} \\
\code{pv/tts.py} & شبیه‌سازی صدا با \en{XTTS-v2} \\
\code{pv/metrics.py} & \en{WER}، \en{CER}، \en{CER} بدون فاصله و شباهت گوینده \\
\code{scripts/prepare\_samples.py} & ساخت نمونه‌های آزمون از \en{FLEURS} \\
\code{scripts/evaluate.py} & ارزیابی کامل و نوشتن \code{results/eval.md} \\
\code{scripts/tune\_clone.py} & مقایسه تنظیمات شبیه‌سازی \\
\code{scripts/rescore.py} & امتیازدهی دوباره با قواعد مختلف نیم‌فاصله \\
\code{scripts/analyze\_stt\_errors.py} & تجزیه خطاها و بازه اطمینان \\
\code{scripts/smoke\_app.py} & آزمون هر دو زبانه رابط از طریق \en{API} \\
\bottomrule
\end{tabularx}
\caption{ساختار کد.}
\end{table}

\subsection{رابط کاربری}

رابط با \en{Gradio} ساخته شده و دو زبانه دارد. در زبانه اول، کاربر فایل صوتی را بارگذاری می‌کند، مدل را از بین \code{fa-turbo} و \code{base-turbo} انتخاب می‌کند و اگر بخواهد متن مرجع هم می‌دهد. متن استخراج‌شده همراه با مدت صوت، زمان پردازش و \en{RTF} نمایش داده می‌شود و در صورت وجود متن مرجع، \en{WER} و \en{CER} هم محاسبه می‌شود. همین صوت به‌طور خودکار نمونه صدای زبانه دوم هم می‌شود تا برای آزمایش سریع لازم نباشد دوباره بارگذاری شود.

در زبانه دوم، نمونه صدا و پاراگراف جدید وارد می‌شود، چک‌باکس رضایت باید تأیید شود و دما را می‌شود بین ۰٫۱ و ۱ تغییر داد (پیش‌فرض ۰٫۵). خروجی، فایل صوتی به‌همراه زمان تولید، \en{RTF} و شباهت گوینده است. کارت مدل نمونه‌ای حدود ۶ تا ۱۵ ثانیه را پیشنهاد می‌کند؛ برای همین اگر طول نمونه خارج از ۶ تا ۳۰ ثانیه باشد، رابط هشدار می‌دهد.

\needspace{10cm}
\section{فناوری‌ها و محیط اجرا}

\subsection{مدل‌ها و داده}

\begin{table}[H]
\centering\small
\begin{tabularx}{\linewidth}{R{2.5cm} L{6.6cm} X}
\toprule
\textbf{کاربرد} & \textbf{منبع} & \textbf{ادعای منبع (تأییدنشده توسط ما)} \\
\midrule
گفتار به متن & \lr{\texttt{MohammadGholizadeh/\allowbreak whisper-large-v3-turbo-persian}} & تنظیم روی حدود ۵۳۴ ساعت گفتار فارسی؛ \en{WER} برابر ۱۹٫۴٪ روی \en{Common Voice 17} فارسی \\
مبنای مقایسه & \lr{\texttt{openai/whisper-large-v3-turbo}} & --- \\
شبیه‌سازی صدا & \lr{\texttt{MohammadJRanjbar/ParsVoice-XTTS}} & آموزش روی ۲٬۲۰۰ ساعت صدا از بیش از ۱٬۸۰۰ گوینده \\
شباهت گوینده & \lr{\texttt{speechbrain/spkrec-ecapa-voxceleb}} & --- \\
داده آزمون & \en{FLEURS}، بخش \en{test} فارسی & --- \\
\bottomrule
\end{tabularx}
\caption{مدل‌ها و داده استفاده‌شده.}
\end{table}

این \en{Whisper} را انتخاب کردم چون طبق کارت مدل علاوه بر گفتار خوانده‌شده، روی یوتیوب و دیالوگ فیلم هم آموزش دیده است. مدل پایه را هم با همان معماری کنارش آزمودم تا مشخص شود تنظیم فارسی چقدر اثر دارد. برای بخش دوم \en{ParsVoice-XTTS} را برداشتم، چون با چند ثانیه صدا کار می‌کند و برای هر گوینده آموزش جداگانه نمی‌خواهد.

\subsection{نرم‌افزار و سخت‌افزار}

\begin{table}[H]
\centering\small
\begin{tabular}{r l r l}
\toprule
\textbf{بسته} & \textbf{نسخه} & \textbf{بسته} & \textbf{نسخه} \\
\midrule
\en{Python} & \en{3.12} & \en{speechbrain} & \en{1.1.2} \\
\en{torch / torchaudio} & \en{2.8.0} & \en{jiwer} & \en{4.0.0} \\
\en{coqui-tts} & \en{0.28.0} & \en{gradio} & \en{6.31.0} \\
\en{transformers} & \en{5.22.1} & \en{soundfile} & \en{0.14.0} \\
\en{accelerate} & \en{1.15.0} & \en{librosa} & \en{1.0.0} \\
\en{huggingface\_hub} & \en{2.4.0} & \en{imageio-ffmpeg} & \en{0.7.0} \\
\bottomrule
\end{tabular}
\caption{نسخه‌های به‌کاررفته (منبع: \code{requirements.txt}).}
\end{table}

آزمون‌ها روی سروری با دو پردازنده \en{Intel Xeon Gold 5220R} (۹۶ رشته)، ۱ ترابایت حافظه و یک کارت \en{NVIDIA RTX~3060} با ۱۲ گیگابایت حافظه انجام شد. آزمون‌های \en{CPU} فقط از ۲۴ رشته استفاده کرده‌اند. از \en{PyTorch 2.8} استفاده شده، چون \en{coqui-tts} از نسخه ۲٫۹ به بعد به \en{torchcodec} و کتابخانه‌های سیستمی \en{FFmpeg} نیاز دارد.

\needspace{7cm}
\section{روش پیاده‌سازی}

\subsection{گفتار به متن}

صوت ورودی با \en{ffmpeg} به ۱۶ کیلوهرتز تک‌کانال تبدیل می‌شود، پس قالب فایل مشکلی ایجاد نمی‌کند. برای رونویسی، بخش‌ها با اندازه دسته ۸ و زبان ثابت فارسی به مدل داده می‌شوند؛ روی \en{GPU} با \en{fp16} و روی \en{CPU} با \en{fp32}.

برای فایل‌های بلند اول همان روش پیش‌فرض کتابخانه را گذاشته بودم، یعنی پنجره‌های همپوشان ۳۰ ثانیه‌ای. وقتی خروجی \code{stt\_03} را با مرجع مقایسه کردم، دیدم عبارت «نیروهای دشمن بوده است. اگر در زمستان از در» درست در مرز پنجره حذف شده است. برای همین برش را خودم نوشتم: صوت بلندتر از ۲۸ ثانیه به قاب‌های ۳۰ میلی‌ثانیه‌ای تقسیم می‌شود و در کم‌انرژی‌ترین قاب بین ثانیه ۱۲ و ۲۸ بریده می‌شود. با این تغییر \en{WER} همان فایل ۲٫۸٪ شد.

\begin{latin}
\begin{codeblock}
def split_on_pauses(wav, max_s=28.0, min_s=12.0):
    hop = int(0.03 * SR)
    segs, start = [], 0
    while len(wav) - start > max_s * SR:
        lo, hi = start + int(min_s * SR), start + int(max_s * SR)
        frames = wav[lo:hi][: (hi - lo) // hop * hop].reshape(-1, hop)
        cut = lo + int(np.argmin((frames ** 2).mean(axis=1))) * hop + hop // 2
        segs.append(wav[start:cut])
        start = cut
    segs.append(wav[start:])
    return segs
\end{codeblock}
\end{latin}

\subsection{شبیه‌سازی صدا}

نمونه صدا پیش از هر کاری از نظر بلندی یکسان می‌شود: صدک ۹۰ از \en{RMS} قاب‌های ۳۰ میلی‌ثانیه‌ای به ۲۳- دسی‌بل می‌رسد و قله سیگنال از ۰٫۹۷ بیشتر نمی‌شود. بردار گوینده را هم با تنظیمات خود مدل در \code{config.json} می‌سازم (\code{gpt\_cond\_len} برابر ۳۰ و \code{gpt\_cond\_chunk\_len} برابر ۴)، چون مدل با همین مقادیر منتشر شده است و نه با مقادیر پیش‌فرض کتابخانه.

متن پیش از تولید کمی آماده می‌شود. حروف «ي» و «ك» یکسان و عددها به حروف نوشته می‌شوند («۱۲» می‌شود «دوازده»)، چون کارت مدل درباره تلفظ نادرست اعداد هشدار داده است. اعشار با «ممیز» و نماد درصد با خود واژه «درصد» جایگزین می‌شود. «ئ»، «ة» و نیم‌فاصله را نگه داشتم، چون در واژگان مدل هستند.

سپس متن در نشانه‌های پایان جمله شکسته می‌شود و جمله‌ای که از ۱۸۰ نویسه بلندتر باشد، در ویرگول یا فاصله بریده می‌شود. هر جمله جداگانه و بدون نشانه پایانی، با دمای ۰٫۵ و جریمه تکرار ۵ تولید می‌شود و جمله‌ها با ۰٫۱۸ ثانیه مکث به هم می‌چسبند. این مکث را دستی گذاشته‌ام و آن را تنظیم نکرده‌ام. خروجی ۲۴ کیلوهرتز است.

یک نکته اجرایی هم داشت: توکنایزر اصلی \en{XTTS} کد زبان \code{fa} را قبول نمی‌کند. در \code{pv/tts.py} این کد را بدون پاک‌سازی اضافه از آن رد کردم، چون متن از قبل آماده شده است.

\subsection{ملاحظه اخلاقی}

تا چک‌باکس رضایت صاحب صدا تأیید نشود، شبیه‌سازی انجام نمی‌شود. این را با \code{scripts/smoke\_app.py} از طریق \en{API} رابط هم امتحان کردم و خطای مربوط در \code{results/app\_test.log} ثبت شده است. البته چک‌باکس فقط اعلام کاربر است و خود سامانه نمی‌تواند رضایت را راستی‌آزمایی کند. نمونه‌های صدای این گزارش از پیکره عمومی \en{FLEURS} برداشته شده‌اند.

\section{روش ارزیابی}\label{sec:method}

\subsection{شاخص‌ها}

\[
\mathrm{WER} = \frac{S + D + I}{N}
\]

در این رابطه $S$ جایگزینی، $D$ حذف، $I$ درج و $N$ تعداد واژه‌های مرجع است. \en{CER} همین نسبت در سطح نویسه است. \en{RTF} زمان پردازش تقسیم بر مدت صوت است و مقدار کمتر از یک یعنی پردازش سریع‌تر از زمان واقعی.

برای \en{CER} رفت‌وبرگشتی، صدای تولیدشده را دوباره با \en{Whisper} رونویسی و با متن ورودی مقایسه کردم. عددها در متن مرجع این مقایسه به حروف نوشته شده‌اند، چون همین‌طور خوانده می‌شوند. این عدد نشان می‌دهد گفتار چقدر قابل فهم است و درباره طبیعی‌بودن آن چیزی نمی‌گوید. برای شباهت گوینده هم شباهت کسینوسی بردارهای \en{ECAPA} نمونه ورودی و خروجی را گرفتم.

برای گفتار به متن، \en{WER} روی کل پیکره محاسبه می‌شود. برای شبیه‌سازی، شاخص هر نمونه جداگانه محاسبه و سپس میانگین گرفته می‌شود.

\subsection{یکسان‌سازی پیش از امتیازدهی}

مرجع و خروجی پیش از مقایسه از یک مسیر می‌گذرند: حروف عربی به فارسی و رقم‌ها به لاتین تبدیل، اعراب و کشیده حذف، عددها به حروف نوشته و نشانه‌گذاری حذف می‌شود. سپس پیشوند «می/نمی» و پسوندهایی مانند «ها/های/تر/ترین» به واژه چسبانده می‌شوند.

\subsection{داده آزمون}

برای گفتار به متن از ۳۰۰ جمله نخست فهرست بخش \en{test} فارسی \en{FLEURS} استفاده کردم (۷۸٫۲ دقیقه و ۷٬۳۱۲ واژه مرجع) که تصادفی انتخاب نشده‌اند. \en{FLEURS} جمله‌های کوتاه دارد، پس برای فایل‌های ۳۰ تا ۶۰ ثانیه‌ای، جمله‌های تصادفی را با ۰٫۴ ثانیه فاصله به هم چسباندم تا طول به حدود ۴۰ ثانیه یا بیشتر برسد (بذر تصادفی ۷). نمونه‌های صدا چهار جمله ۸ تا ۱۵ ثانیه‌ای از همین پیکره‌اند. متن‌های جدید سه پاراگراف زیر هستند که در هیچ‌یک از نمونه‌ها نیامده‌اند:

\begin{enumerate}
\item پروژه ساختمان مسکونی در ۱۲ طبقه طراحی شده و عملیات اجرایی آن از بهار سال آینده آغاز می‌شود. برآورد اولیه نشان می‌دهد که حدود ۸۵۰۰ متر مکعب بتن برای اسکلت لازم است.
\item صبح امروز هوای تهران نسبتاً پاک بود و بسیاری از مردم برای پیاده‌روی به پارک‌ها رفتند. کارشناسان هواشناسی پیش‌بینی می‌کنند که از فردا دما کاهش پیدا کند.
\item هوش مصنوعی در سال‌های اخیر پیشرفت چشمگیری داشته است. امروزه مدل‌های زبانی می‌توانند متن بنویسند، ترجمه کنند و حتی به پرسش‌های تخصصی پاسخ دهند.
\end{enumerate}

\clearpage
\section{نتایج گفتار به متن}\label{sec:stt}

\subsection{نتیجه کلی}

\begin{table}[H]
\centering\small
\begin{tabular}{r c c c c c}
\toprule
\textbf{مدل} & \en{WER} & \en{CER} & \textbf{\en{CER} بدون فاصله} & \en{RTF} & \textbf{حافظه \en{GPU}} \\
\midrule
\en{Whisper turbo} فارسی & \textbf{۸٫۹٪} & ۲٫۸٪ & ۲٫۸٪ & ۰٫۰۴۴ & ۲٫۱۱ \en{GB} \\
\en{Whisper turbo} پایه & ۱۸٫۰٪ & ۵٫۴٪ & ۵٫۷٪ & ۰٫۰۴۴ & ۲٫۱۱ \en{GB} \\
\bottomrule
\end{tabular}
\caption{نتایج روی ۳۰۰ جمله با \en{RTX~3060} (منبع: \code{results/eval.md}).}
\end{table}

\begin{table}[H]
\centering\small
\begin{tabular}{r c c c c c}
\toprule
\textbf{فایل} & \textbf{طول} & \multicolumn{2}{c}{\textbf{مدل فارسی}} & \multicolumn{2}{c}{\textbf{مدل پایه}} \\
 & & \en{WER} & \en{CER} & \en{WER} & \en{CER} \\
\midrule
\code{stt\_01.wav} & ۴۳٫۰ ثانیه & ۳٫۳٪ & ۰٫۴٪ & ۱۸٫۳٪ & ۴٫۱٪ \\
\code{stt\_02.wav} & ۵۱٫۷ ثانیه & ۷٫۶٪ & ۲٫۲٪ & ۱۶٫۵٪ & ۴٫۳٪ \\
\code{stt\_03.wav} & ۴۰٫۸ ثانیه & ۲٫۸٪ & ۱٫۱٪ & ۵٫۶٪ & ۱٫۹٪ \\
\bottomrule
\end{tabular}
\caption{فایل‌های ۴۱ تا ۵۲ ثانیه‌ای. مدل فارسی هر فایل را در ۱٫۵ تا ۱٫۹ ثانیه پردازش کرد.}
\end{table}

\subsection{تجزیه خطاها}

\begin{table}[H]
\centering\small
\begin{tabular}{r c c}
\toprule
\textbf{مورد} & \textbf{مدل فارسی} & \textbf{مدل پایه} \\
\midrule
واژه‌های درست از ۷٬۳۱۲ & ۶٬۸۰۱ & ۶٬۱۵۸ \\
جایگزینی & ۴۴۰ & ۱٬۰۴۶ \\
حذف & ۷۱ & ۱۰۸ \\
درج & ۱۳۸ & ۱۵۹ \\
\textbf{\en{WER}} & \textbf{۸٫۸۸٪} & \textbf{۱۷٫۹۶٪} \\
بازه اطمینان ۹۵٪ & ۷٫۸۷٪ تا ۹٫۹۵٪ & ۱۶٫۵۴٪ تا ۱۹٫۴۰٪ \\
جمله‌های بدون خطا از ۳۰۰ & ۷۱ & ۱۸ \\
میانه \en{WER} جمله‌ها & ۶٫۶٪ & ۱۶٫۷٪ \\
صدک ۹۰ \en{WER} جمله‌ها & ۲۰٫۹٪ & ۳۵٫۰٪ \\
\bottomrule
\end{tabular}
\caption{تجزیه خطاها (منبع: \code{results/stt\_errors.md}).}
\end{table}

بازه‌های اطمینان را با ۲۰۰۰ بار بازنمونه‌گیری در سطح جمله حساب کردم. اختلاف \en{WER} دو مدل روی همان جمله‌ها بین ۸٫۰۱٪ تا ۱۰٫۱۶٪ به نفع مدل فارسی است و بیشتر این بهبود از کمترشدن جایگزینی‌ها می‌آید (از ۱٬۰۴۶ به ۴۴۰).

\begin{table}[H]
\centering\small
\begin{tabular}{r r c}
\toprule
\textbf{نوع خطا} & \textbf{نمونه‌ها (مرجع ← خروجی)} & \textbf{تعداد} \\
\midrule
واژه‌های هم‌آوا & اتو ← اطو، سهره ← سحره، خصایص ← خصایث، غرش ← قرش، حرا ← هرا & ۲۲ \\
نگارش کسره اضافه & همه ← همهی، دهه ← دههی، دوره ← دورهی & ۱۵ \\
مطابقت فعل & کنند ← کند، میشوند ← میشود & ۸ \\
شکل کوتاه فعل & ست ← است & ۵ \\
گ و ک & وبلاگها ← وبلاکها & ۵ \\
\bottomrule
\end{tabular}
\caption{دوازده جایگزینی پرتکرار مدل فارسی، گروه‌بندی‌شده. مجموع ۵۵ مورد از ۴۴۰ جایگزینی.}
\end{table}

وقتی این جایگزینی‌ها را نگاه کردم، دیدم بیشترشان اشتباه شنیداری نیستند؛ مدل صدا را درست گرفته ولی املای دیگری نوشته است. در فارسی چند حرف هم‌آوا داریم و وقتی بافت جمله کمکی نکند، انتخاب بینشان سخت است.

\subsection{اثر قاعده امتیازدهی}

\en{WER} فارسی خیلی به نحوه شمردن فاصله‌ها وابسته است؛ «می‌روم»، «می روم» و «میروم» یک واژه‌اند. قاعده نهایی را بعد از دیدن نتایج انتخاب کردم، چون دو قاعده ساده‌تر به نفع یکی از مدل‌ها تمام می‌شدند. برای اینکه این انتخاب پنهان نماند، نتیجه هر سه را آورده‌ام.

\begin{table}[H]
\centering\small
\begin{tabular}{r c c}
\toprule
\textbf{قاعده نیم‌فاصله} & \textbf{\en{WER} مدل فارسی} & \textbf{\en{WER} مدل پایه} \\
\midrule
تبدیل به فاصله & ۱۳٫۹۸٪ & ۱۸٫۸۴٪ \\
حذف و چسباندن & ۱۶٫۶۵٪ & ۲۷٫۹۴٪ \\
\textbf{چسباندن پیشوند و پسوند (نهایی)} & \textbf{۸٫۸۸٪} & \textbf{۱۷٫۹۶٪} \\
\bottomrule
\end{tabular}
\caption{اثر قاعده امتیازدهی (منبع: \code{results/rescore.txt}).}
\end{table}

\en{CER} بدون فاصله به این انتخاب وابسته نیست و مقایسه‌ای محافظه‌کارانه‌تر است: ۲٫۸٪ در برابر ۵٫۷٪. عدد ۸٫۹٪ را نباید مستقیم با ۱۹٫۴٪ کارت مدل مقایسه کرد، چون داده و قاعده امتیازدهی متفاوت‌اند.

\subsection{نمونه‌های ورودی و خروجی}

فایل‌های صوتی و متن‌های زیر در پوشه \code{examples/part1\_stt} قرار دارند. واژه‌های پررنگ خطای واقعی‌اند؛ تفاوت‌های نیم‌فاصله و نشانه‌گذاری در امتیازدهی حذف می‌شوند.

\textbf{نمونه ۱ (\code{stt\_03}، مرجع):} «نیازی به گفتن نیست، اگر یکی از زبان‌های رومی را بلد باشید، یادگیری زبان پرتغالی برای شما آسان‌تر خواهد بود. برخورد در بالای زمین‌های کوهستانی رخ داد و اعتقاد می‌رود که در نتیجه شلیک نیروهای دشمن بوده است. اگر در زمستان از دریای بالتیک شمالی عبور می‌کنید، محل کابین را بررسی کنید زیرا هنگام عبور از یخ برای کسانی که بیشتر تحت تأثیر قرار گرفته باشند صداهای نسبتاً وحشتناکی شنیده می‌شود.»

\textbf{خروجی:} «نیازی به گفتن نیست. اگر یکی از زبانهای رومی را بلد باشید، یادگیری زبان پرتغالی برای شما آسانتر خواهد بود. برخورد در بالای زمینهای کوهستانی رخ داد و اعتقاد میرود که در نتیجه شلیک نیروهای دشمن بوده است. اگر در زمستان از دریای بالتیک شمالی عبور میکنید محل کابین را بررسی کنید زیرا هنگام عبور از یخ برای کسانی که بیشتر تحت \textbf{تحصیل} قرار گرفته باشند \textbf{صداهایی} نسبتا وحشتناکی شنیده میشود»

\textbf{نمونه ۲ (\code{stt\_01}، مرجع):} «در شمال، این منطقه توسط ساحل صحرا و از جنوب و غرب توسط اقیانوس اطلس احاطه شده است. دو پودر خشک را مخلوط کنید و سپس با دستان خشک و تمیز آنها را فشار دهید تا به‌ توپی در آیند. عکاس‌ها سپس جای یک زن سالخورده که نیاز به دستشویی داشت را گرفتند. مندوزا با گلوله از پای درآورده شد.»

\textbf{خروجی:} «در شمال، این منطقه توسط ساحل صحرا و از جنوب و غرب توسط اقیانوس اطلس احاطه شده است. دو پودر خشک را مخلوط کنید و سپس با دستان خشک و تمیز آنها را فشار دهید تا به توپی درآیند. عکاسها سپس جای یک زن سالخورده که نیاز به دستشویی داشت را گرفتند، مندوزا با گلوله از پای درآورده شد.»

\textbf{نمونه ۳ (\code{sentence\_106}، مرجع):} «اگر فقط توسط کشتی‌های تفریحی در ساحل پهلو بگیرید، نیازی به ویزای مجزا نخواهید داشت (از سال 2009).»

\textbf{خروجی:} «اگر فقط توسط کشتی های تفریحی در ساحل پهلو بگیرید \textbf{نیاز} به ویزای \textbf{مجازه} نخواهید داشت از سال ۲۰۰۹»

\clearpage
\section{نتایج شبیه‌سازی صدا}\label{sec:tts}

\subsection{اثر هر اصلاح}

آزمون تنظیم روی ۳ گوینده، ۲ متن و ۲ بذر تصادفی، یعنی ۱۲ نمونه برای هر حالت انجام شد. حالت الف با دمای ۰٫۷، بدون یکسان‌سازی بلندی و با تنظیمات پیش‌فرض کتابخانه برای بردار گوینده اجرا شد.

\begin{figure}[H]
\centering
\begin{latin}
\begin{tikzpicture}[x=0.42cm,y=0.72cm]
  \foreach \y/\v/\t/\f in {4/22.7/{۲۲٫۷٪}/{الف: پیش از اصلاحات},3/10.2/{۱۰٫۲٪}/{ب: بردار گوینده و بلندی},2/10.1/{۱۰٫۱٪}/{ج: ب + دمای ۰٫۵},1/5.0/{۵٫۰٪}/{د: ب + حذف نشانه پایانی}} {
    \fill[accent!55,rounded corners=1pt] (-\v,\y-0.2) rectangle (0,\y+0.2);
    \node[anchor=east,font=\small] at (-\v-0.2,\y) {\rl{\t}};
    \node[anchor=west,font=\small] at (0.3,\y) {\rl{\f}};
  }
  \fill[accent,rounded corners=1pt] (-3.6,-0.2) rectangle (0,0.2);
  \node[anchor=east,font=\small\bfseries] at (-3.8,0) {\rl{۳٫۶٪}};
  \node[anchor=west,font=\small\bfseries] at (0.3,0) {\rl{ه: ب + دمای ۰٫۵ + حذف نشانه}};
  \draw[gray!60] (0,-0.5) -- (0,4.5);
\end{tikzpicture}
\end{latin}
\caption{میانگین \en{CER} رفت‌وبرگشتی هر حالت. حالت نهایی پررنگ‌تر است (منبع: \code{results/tune.md}).}
\end{figure}

\begin{table}[H]
\centering\small
\begin{tabular}{r c c c c c}
\toprule
\textbf{حالت} & \en{WER} & \en{CER} & \textbf{شباهت} & \textbf{طول صوت} \\
\midrule
الف: پیش از اصلاحات & ۳۶٫۵٪ & ۲۲٫۷٪ & ۰٫۶۵ & ۲۰٫۳ ثانیه \\
ب: تنظیمات بردار گوینده و یکسان‌سازی بلندی & ۲۰٫۵٪ & ۱۰٫۲٪ & ۰٫۶۷ & ۱۵٫۵ ثانیه \\
ج: ب + دمای ۰٫۵ & ۲۰٫۱٪ & ۱۰٫۱٪ & ۰٫۶۷ & ۱۵٫۹ ثانیه \\
د: ب + حذف نشانه پایانی & ۱۵٫۴٪ & ۵٫۰٪ & ۰٫۶۶ & ۱۴٫۸ ثانیه \\
\textbf{ه: ب + دمای ۰٫۵ + حذف نشانه پایانی} & \textbf{۹٫۵٪} & \textbf{۳٫۶٪} & \textbf{۰٫۶۹} & \textbf{۱۶٫۵ ثانیه} \\
\bottomrule
\end{tabular}
\caption{مقایسه تنظیمات.}
\end{table}

\begin{table}[H]
\centering\small
\begin{tabular}{r c c c c c c}
\toprule
\textbf{گوینده} & \multicolumn{3}{c}{\textbf{الف: پیش از اصلاحات}} & \multicolumn{3}{c}{\textbf{ه: نهایی}} \\
 & \en{WER} & \en{CER} & \textbf{شباهت} & \en{WER} & \en{CER} & \textbf{شباهت} \\
\midrule
\code{voice\_01} & ۲۱٫۵٪ & ۱۰٫۵٪ & ۰٫۷۳ & ۱۰٫۶٪ & ۴٫۲٪ & ۰٫۷۱ \\
\code{voice\_03} & ۶۲٫۰٪ & ۴۴٫۹٪ & ۰٫۵۸ & ۵٫۴٪ & ۱٫۲٪ & ۰٫۷۱ \\
\code{voice\_04} & ۲۵٫۹٪ & ۱۲٫۸٪ & ۰٫۶۵ & ۱۲٫۵٪ & ۵٫۴٪ & ۰٫۶۴ \\
\bottomrule
\end{tabular}
\caption{اثر اصلاحات به تفکیک گوینده.}
\end{table}

نتیجه اجرای اول خوب نبود و دو مشکل عمده داشت. اولی به \code{voice\_03} برمی‌گشت که سطح صدایش حدود ۵۵- دسی‌بل بود؛ \en{CER} خروجی‌های این گوینده بین ۳۸٪ تا ۱۷۸٪ شد. بعد از یکسان‌سازی بلندی، میانگین \en{CER} همین گوینده از ۴۴٫۹٪ به ۷٫۴٪ آمد (حالت الف به ب).

مشکل دوم این بود که بعد از پایان جمله‌ها هجاهای اضافه تولید می‌شد و \en{Whisper} چیزهایی مثل «پانتقی» می‌شنید. بیشترین اثر را حذف نشانه پایانی داشت. هنوز نمی‌دانم چرا این کار جواب می‌دهد؛ علتش را بررسی نکردم و فقط اثرش را اندازه گرفتم.

\subsection{نتیجه نهایی}

ارزیابی نهایی روی ۴ گوینده و ۳ پاراگراف، با \en{CPU} و ۲۴ رشته انجام شد. نمونه‌های «کنارگذاشته‌شده» شامل \code{voice\_02} یا پاراگراف ۳ هستند که در انتخاب تنظیمات به کار نرفته‌اند.

\begin{table}[H]
\centering\small
\begin{tabular}{r c c c c}
\toprule
\textbf{حالت} & \en{WER} & \en{CER} & \textbf{شباهت} & \en{RTF} \\
\midrule
پیش از اصلاحات (روی \en{GPU}) & ۳۹٫۵٪ & ۲۹٫۰٪ & ۰٫۶۵ & ۰٫۳۲ \\
\textbf{نهایی، یک بار تولید} & \textbf{۱۰٫۶٪} & \textbf{۳٫۴٪} & \textbf{۰٫۶۹} & \textbf{۱٫۰۷} \\
نهایی، فقط ۶ نمونه کنارگذاشته‌شده & ۱۳٫۶٪ & ۴٫۰٪ & ۰٫۷۰ & --- \\
نهایی، تا ۳ بار تولید برای هر جمله & ۱۲٫۷٪ & ۴٫۷٪ & ۰٫۶۹ & ۳٫۸۶ \\
\bottomrule
\end{tabular}
\caption{نتایج نهایی (منبع: \code{results/eval.md}).}
\end{table}

\begin{table}[H]
\centering\small
\begin{tabular}{r c c c c c}
\toprule
\textbf{گوینده} & \textbf{پاراگراف} & \textbf{طول صوت} & \en{RTF} & \textbf{شباهت} & \textbf{\en{CER} رفت‌وبرگشتی} \\
\midrule
\code{voice\_01} & ۱ & ۱۶٫۶۲ ثانیه & ۱٫۱۰ & ۰٫۷۱ & ۵٫۱٪ \\
\code{voice\_01} & ۲ & ۱۴٫۶۱ ثانیه & ۱٫۱۴ & ۰٫۷۰ & ۲٫۷٪ \\
\code{voice\_01} & ۳ & ۱۵٫۷۷ ثانیه & ۱٫۱۷ & ۰٫۷۱ & ۳٫۷٪ \\
\code{voice\_02} & ۱ & ۱۶٫۳۷ ثانیه & ۱٫۰۵ & ۰٫۷۳ & ۵٫۷٪ \\
\code{voice\_02} & ۲ & ۱۲٫۴۸ ثانیه & ۱٫۰۶ & ۰٫۷۶ & ۲٫۷٪ \\
\code{voice\_02} & ۳ & ۱۲٫۸۹ ثانیه & ۱٫۰۹ & ۰٫۷۳ & ۳٫۰٪ \\
\code{voice\_03} & ۱ & ۱۶٫۰۹ ثانیه & ۱٫۰۶ & ۰٫۷۱ & ۰٫۶٪ \\
\code{voice\_03} & ۲ & ۱۱٫۷۹ ثانیه & ۱٫۰۴ & ۰٫۷۰ & ۰٫۷٪ \\
\code{voice\_03} & ۳ & ۱۱٫۶۹ ثانیه & ۱٫۰۴ & ۰٫۶۵ & ۵٫۹٪ \\
\code{voice\_04} & ۱ & ۲۰٫۶۹ ثانیه & ۱٫۰۳ & ۰٫۶۳ & ۵٫۷٪ \\
\code{voice\_04} & ۲ & ۱۵٫۴۴ ثانیه & ۱٫۰۱ & ۰٫۶۵ & ۱٫۴٪ \\
\code{voice\_04} & ۳ & ۱۳٫۲۳ ثانیه & ۱٫۱۰ & ۰٫۶۶ & ۳٫۰٪ \\
\bottomrule
\end{tabular}
\caption{جزئیات هر نمونه با تنظیمات نهایی، یک بار تولید.}
\end{table}

برای مقایسه، میانگین شباهت میان خود نمونه‌های مرجع ۰٫۲۸ است. \en{FLEURS} شناسه گوینده ندارد، پس ممکن است برخی از آن‌ها یک نفر باشند و این عدد فقط یک مبنای تقریبی است.

یک ایده هم جواب نداد: هر جمله را تا سه بار تولید کردم و بهترین نسخه را با \en{Whisper} برداشتم. سرعت ۳٫۶ برابر کمتر شد و نتیجه بهتر نشد، پس آن را از رابط برداشتم. البته با ۱۲ نمونه نمی‌شود قطعی نتیجه گرفت.

\subsection{نمونه‌های ورودی و خروجی}

برای هر گوینده، نمونه مرجع و سه صدای تولیدشده در پوشه \code{examples/part2\_cloning} قرار دارد و جدول کامل آن در \code{results.tsv} همان پوشه است. متن‌های زیر آن چیزی‌اند که \en{Whisper} از سه خروجی حالت «ه» شنیده است؛ بنابراین بخشی از خطاها ممکن است از خود \en{Whisper} باشد.

\textbf{بهترین نمونه} (\code{voice\_03}، پاراگراف ۱، \en{CER} برابر ۰٫۶٪): «پروژه ساختمان مسکونی در \textbf{دوازد} طبقه طراحی شده و عملیات اجرایی آن از بهار سال آینده آغاز می شود برآورد اولیه نشان می دهد که حدود هشت هزار و پانصد متر مکعب بتن برای اسکلت لازم است»

\textbf{نمونه ضعیف} (\code{voice\_01}، پاراگراف ۱، \en{CER} برابر ۶٫۳٪): «\textbf{هوژه ی} ساختمان مسکونی در دوازده طبقه طراحی شده و عملیات \textbf{اجرای} آن از بهار سال آینده آغاز می شود برآورد اولیه نشان می دهد که حدود هشت هزار و پانصد متر \textbf{مکعبتاً برای یه} اسکلت لازم است»

\textbf{بدترین نمونه} (\code{voice\_04}، پاراگراف ۲، \en{CER} برابر ۱۱٫۶٪، ۲۷٫۱ ثانیه): «صبح امروز، هوای تهران نسبتا پاک بود و بسیاری از مردم برای پیاده روی به پارک ها \textbf{رفتندن} کارشناسان هواشناسی پیش بینی \textbf{میکنن} که از فردا... \textbf{خیلی خوش اومدیم...}». مدل در این نمونه انتهای جمله را ناتمام گذاشته و عبارتی ساختگی تولید کرده است.

\section{سرعت و منابع}\label{sec:speed}

\begin{table}[H]
\centering\small
\begin{tabular}{r c r}
\toprule
\textbf{مورد} & \en{RTF} & \textbf{شرایط} \\
\midrule
گفتار به متن & ۰٫۰۴۴ & \en{GPU}، با ۲٫۱۱ \en{GB} حافظه \\
گفتار به متن & ۰٫۵۵ & \en{CPU} با ۲۴ رشته، فقط یک فایل (\code{stt\_01.wav}) \\
شبیه‌سازی صدا & ۱٫۰۷ & \en{CPU} با ۲۴ رشته، تنظیمات نهایی \\
شبیه‌سازی صدا & ۳٫۸۶ & \en{CPU} با ۲۴ رشته، تا ۳ بار تولید \\
شبیه‌سازی صدا & ۰٫۳۲ & \en{GPU} با ۳٫۶۷ \en{GB} حافظه، \textbf{پیش از اصلاحات} \\
\bottomrule
\end{tabular}
\caption{سرعت و حافظه.}
\end{table}

با این ضریب‌ها، رونویسی یک فایل ۶۰ ثانیه‌ای حدود ۲٫۶ ثانیه روی \en{GPU} و حدود ۳۳ ثانیه روی \en{CPU} طول می‌کشد، و تولید ۱۵ ثانیه گفتار روی \en{CPU} حدود ۱۶ ثانیه. این سه عدد از ضریب‌های بالا محاسبه شده‌اند و جداگانه اندازه‌گیری نشده‌اند. بارگذاری مدل فارسی گفتار به متن ۳٫۱ ثانیه طول کشید.

سرعت تنظیمات نهایی روی \en{GPU} اندازه‌گیری نشده است، چون \en{GPU} در این مدت در اختیار کار دیگری بود.

\clearpage
\section{راهنمای نصب و اجرا}\label{sec:guide}

پیش‌نیازها لینوکس و \en{Python 3.12} است. راهنمای کامل در \code{README.md} آمده است.

\begin{latin}
\begin{codeblock}
python3.12 -m venv .venv
.venv/bin/pip install torch==2.8.0 torchaudio==2.8.0 \
    --index-url https://download.pytorch.org/whl/cu128
.venv/bin/pip install -r requirements.txt

# ParsVoice-XTTS is gated: accept its terms on the model page first
.venv/bin/hf auth login

.venv/bin/python app.py                        # http://localhost:7860
CUDA_VISIBLE_DEVICES= .venv/bin/python app.py  # CPU only
\end{codeblock}
\end{latin}

\textbf{بازتولید نتایج:}

\begin{latin}
\begin{codeblock}
.venv/bin/python -m scripts.prepare_samples
.venv/bin/python -m scripts.evaluate --stt-corpus 300 --tts
.venv/bin/python -m scripts.tune_clone
.venv/bin/python -m scripts.rescore
.venv/bin/python -m scripts.analyze_stt_errors
\end{codeblock}
\end{latin}

\section{محدودیت‌ها}\label{sec:limits}

مهم‌ترین کمبود این کار دو چیز است. اول اینکه طبیعی‌بودن صدا را با هیچ شاخصی نسنجیده‌ام، در حالی که از معیارهای ارزیابی است و باید با شنیدن فایل‌ها قضاوت شود. حذف نشانه پایانی هم ممکن است لحن پایان جمله را عوض کرده باشد. دوم گفتار محاوره‌ای است: برایش راهکار خاصی نساختم و آزمونی هم انجام ندادم. تنها کاری که در این جهت شد، انتخاب مدلی است که طبق کارتش روی یوتیوب و دیالوگ فیلم هم آموزش دیده است. \en{FLEURS} گفتار خوانده‌شده است و کارت مدل برای دیالوگ فیلم (مجموعه \en{Filimo}) \en{WER} برابر ۳۱٫۲٪ گزارش کرده، پس روی گفتار محاوره‌ای باید انتظار دقت کمتری داشت. بخش تولید گفتار هم شکل‌های محاوره‌ای را جداگانه پردازش نمی‌کند.

داده آزمون هم محدود است. همه نمونه‌های بخش \en{test} فارسی \en{FLEURS} برچسب مرد دارند، پس صدای زنانه آزموده نشده است. نمونه‌های صدا هم تمیزند و رفتار سامانه با صدای نویزی را نمی‌دانم. نتایج شبیه‌سازی فقط روی ۱۲ نمونه است و در ارزیابی نهایی یک بذر تصادفی به کار رفته. خطا هنوز رخ می‌دهد؛ مثلاً یکی از ۱۲ نمونه حالت «ه» به ۲۷ ثانیه رسید و عبارتی ساختگی تولید کرد. ۳۰۰ جمله گفتار به متن تصادفی انتخاب نشده‌اند و فایل‌های ۳۰ تا ۶۰ ثانیه‌ای هم از چسباندن جمله‌ها ساخته شده‌اند و گفتار پیوسته نیستند. دقت برای طول‌های بیشتر در این گزارش نیامده است.

درباره شاخص‌ها هم باید احتیاط کرد. خطای \en{CER} رفت‌وبرگشتی می‌تواند از خود \en{Whisper} باشد و همین مدل در انتخاب تنظیمات هم به کار رفته است. شاخص شباهت گوینده کالیبره نشده است؛ \en{ECAPA} روی \en{VoxCeleb} آموزش دیده و آستانه‌ای برای «یک گوینده بودن» در فارسی ندارم. قاعده امتیازدهی را هم بعد از دیدن نتایج انتخاب کردم و بازه‌های اطمینان فقط تغییرپذیری جمله‌ها را نشان می‌دهند.

\section{پیشنهادهای توسعه}

اگر فرصت بیشتری داشتم، اول سراغ گفتار محاوره‌ای می‌رفتم: داده محاوره‌ای با متن مرجع جمع می‌کردم، قاعده‌ای می‌نوشتم که «میرم» و «می‌روم» را در امتیازدهی منصفانه مقایسه کند و بعد مدل را روی همین داده و ضبط‌های واقعی با نویز، با رضایت صاحبان صدا، بیشتر تنظیم می‌کردم. قدم بعدی یک آزمون شنیداری با چند شنونده و گزارش \en{MOS} است تا طبیعی‌بودن صدا هم عدد داشته باشد.

برای خطاهای املایی بخش~\ref{sec:stt}، یک پس‌پردازش با مدل زبانی می‌تواند واژه‌های هم‌آوا و نگارش کسره اضافه را اصلاح کند. در بخش شبیه‌سازی هم باید صدای زنانه آزمون شود، شاخص شباهت با جفت‌های یک گوینده و دو گوینده در فارسی کالیبره شود و علت هجاهای اضافه پیدا شود؛ بریدن خودکار انتهای صدا با زمان‌بندی واژه‌ها هم راه‌حل احتمالی است. سنجش تنظیمات نهایی روی \en{GPU} و مقایسه با مدل‌های دیگری مثل \en{Chatterbox} فارسی هم کارهای کوچک‌تری هستند که انجام نشده‌اند.

\end{document}
